% Full GPSS separation construction, guarded comparison, and real countermodel.

\section{The past/future separation construction}
\label{sec:annex-p09-gpss-past-future-separation}

\begin{annexitem}{annex:p9.gpss.separation.translation}{thm:p09-gpss-separation-translation}{Every temporal constructor and its exact first-order image}
\begin{proof}
Fix a nonempty strict linear-order model and a point \(t\).
An atomic proposition holds at \(t\) exactly when the same
point belongs to its unary interpretation, which is the
atomic first-order clause. Constants and Boolean operations
have their ordinary truth conditions in both languages.

Suppose the assertion is known for \(A,B\).
The until formula holds at \(t\) precisely when there
is \(s>t\) at which \(B\) holds and \(A\) holds at
every \(u\) with \(t<u<s\). Apply the induction
hypothesis separately at \(s\) and at each such \(u\).
The result is exactly its displayed first-order formula
with \(x=t\) and witness \(y=s\). Conversely, the
first-order witness and its universal matrix give
those same temporal truths by the induction hypothesis.

For since, its witness instead satisfies \(s<t\).
The interval condition concerns \(s<u<t\).
Applying the same induction hypotheses gives its
displayed first-order clause, in both directions.
In neither clause is the interval operand required
at either external endpoint. If the open interval
is empty its universal condition is true, but the
strict endpoint witness is still required.

All copied bound names are chosen fresh, so the
free current variable receives the intended witness
or intervening point and is not captured by another
constituent. Induction on the finite syntax proves
both the equivalence and termination of the translation.
Unary predicates are copied, never newly introduced.

The source operand order is preserved: the infix
first argument labels the interval and the second
the witness. Kamp's prefix notation instead lists
the witness first, giving precisely the two
displayed correspondences.

The derived \(F,P\) clauses use a true interval
operand and therefore require only a strict future
or past witness. Negating the corresponding
existential failure gives \(G,H\). A false
interval operand requires absence of an intervening
point; thus \(X,Y\) are exactly the conditional
successor and predecessor operations. These
calculations are valid on every linear order.
They include a missing successor or predecessor,
where the respective formula is false.

Finally, prove locality directly by syntax induction
on each pure fragment. A future-until formula at
\(t\) only uses points in the tail \([t,+\infty)\);
its operand truth there is unchanged when that
tail is preserved, by induction. A pure-since
formula uses only \((-\infty,t]\), and the
same induction holds with earlier witnesses.
Atoms may still test \(t\). This proves the
two precise locality statements, without an
order-completeness or endpoint assumption.
\end{proof}
\end{annexitem}

\begin{annexitem}{annex:p9.gpss.separation.dual}{thm:p09-gpss-separation-dual}{The formula, model and regional reversal identities}
\begin{proof}
For an order atom \(u<v\), its reversed formula
holds in \(\mathcal M\) exactly when \(v<u\)
holds there, which is exactly the original atom
in \(\mathcal M^{\mathrm{op}}\). Equality and
unary atoms have unchanged interpretations.
Boolean operations preserve this identity.
Quantification ranges over the same domain on
both sides, so induction on first-order syntax
proves the formula identity for every assignment.

For temporal syntax, take an until formula.
In the opposite model its witness \(s\) satisfies
\(t<_{\mathrm{op}}s\), meaning \(s<t\) in
the original. Its intervening points satisfy
\(s<u<t\) there. By induction on the two
operands, this is exactly
\(A^{\mathrm d}\mathbin{\mathsf S_s}B^{\mathrm d}\)
in the original model. The since case gives
until in the same way. Neither calculation
exchanges operand positions. Atoms, constants
and Boolean operations complete the induction.
Repeating reversal restores every constructor.
The first-order translation has exactly the
same witness and interval reversals, so its
commutation identity follows as well, with
only fresh bound-variable names differing.

Reversal preserves strict linearity: asymmetry,
transitivity and comparison of distinct points
are preserved when all inequalities are reversed.
A successor in the opposite order is a
predecessor in the original, and vice versa.
The condition that such a neighbor exists
whenever there is a point in that direction
therefore preserves discreteness. It also
exchanges the first and last points rather
than requiring them to be absent.

An upper bound in the opposite order is a
lower bound in the original. Its least
upper bound is precisely the original
greatest lower bound. Reversing the
other direction exchanges infima and
suprema as well. The two defining
clauses of Dedekind completeness
therefore hold in the opposite
order whenever they hold in the
original. Nonemptiness of the
domain is unchanged.

For the regional assertion use
induction on the exact GPSS stage.
At level zero both classes are
empty. At the next level a
positive component has one
strictly internal witness and
finitely many earlier regional
or point children. Reversal
takes the guard \(a<z<b\)
to \(b<z<a\), after exchanging
the displayed conjunction order.
It takes a left child on
\((a,z)\) to a right child
on \((z,a)\), and a right
child on \((z,b)\) to a
left child on \((b,z)\).
A signed unary literal at
\(z\) remains that literal.
Thus the reversed component
belongs to the stated next
stage on the reversed
endpoint pair. Its negative
sign is retained by the
Boolean reversal. Applying
the same argument twice
proves equality of the
two translated classes
up to guard normalization.

For a before component
the omitted lower bound
leaves only \(z<x\);
reversal gives \(x<z\)
with its upper bound
omitted. This is the
after constructor, and
its children reverse
as just proved. The
formal \(-\infty\)
and \(+\infty\)
are descriptions of
these omitted guards,
not elements of either
model.

All quantified points
of an after component
lie strictly after its
outer current point;
the nested left and
right regions remain
inside that tail.
Induction through the
same grammar proves
future locality.
Before components
have the corresponding
past locality, either
by this induction
or by reversal.
At a first point
a positive outer
before existential
has no witness,
and its negation
is true. In the
opposite order these
are exactly the
after values at
a last point.
This verifies the
empty-ray cases
within the same
typed construction.
\end{proof}
\end{annexitem}

\begin{annexitem}{annex:p9.gpss.past.future.separation}{thm:p09-gpss-past-future-separation}{The full effective separation at every discrete-complete point}
\begin{proof}
Fix a discrete complete model
\(\mathcal M\) and any point \(t\).
The temporal-to-first-order theorem
identifies \(A\) at that point
with \(A^*(x)\) under \(x=t\).
Apply the fully proved GPSS
special normal form with the
declared one-variable list \((x)\).
There is only one possible weak
order type on that list; its
prefix \(x=x\) is true. Thus
every disjunct has exactly the
before, current and after
grouping used by the algorithm.
The normal-form theorem supplies
all quantifier and equality
cases before this grouping;
no further semantic choice
is made here.

For an after component
\(\gamma(x)\), the completed
discrete future construction
gives \(T_F^{\mathrm{dc}}(\gamma)\)
in \(L(U)\), equivalent at
every point of every
discrete complete model.
Its hypothesis is satisfied
because an after component
is future by the regional
lemma. There is no restriction
to the first point.

For a before component
\(\beta(x)\), its first-order
reverse is an after component.
The opposite model is also
discrete and complete.
Consequently the exact
four-step identity is
\[
\begin{aligned}
 \mathcal M,t\models\beta
 &\Longleftrightarrow
       \mathcal M^{\mathrm{op}},t\models\beta^{\mathrm d}\\
 &\Longleftrightarrow
       \mathcal M^{\mathrm{op}},t\models
                    T_F^{\mathrm{dc}}(\beta^{\mathrm d})\\
 &\Longleftrightarrow
       \mathcal M,t\models
          \bigl(T_F^{\mathrm{dc}}(\beta^{\mathrm d})\bigr)^{\mathrm d}.
\end{aligned}
\]
The final formula uses only since,
because the preceding temporal
formula used only until.
This proves the past translator
with its complete typing and
scope. It does not attempt to
complement a future-labelled
formula on a formal left ray
in the original model.

Replace each regional conjunct
by the corresponding translation
in each clause. The current
signed literals have exactly
their temporal atomic truth
values. Their Boolean conjunction
belongs to both pure fragments.
Finite conjunction and disjunction
preserve the established
equivalences, giving precisely
\(\operatorname{Sep}_{\mathrm{dc}}(A)\).
The reverse implication follows
from the reverse direction of
each same equivalence and
the normal-form equivalence.
Thus the algorithm proves
the original formula and
its separated output equivalent,
not merely equisatisfiable.

Every step terminates. The
input syntax and its
first-order translation are
finite. The normal-form
procedure returns a finite
list. Each regional translation
is the supplied finite
construction, and reversal
is a finite traversal.
The output combines only
finitely many such formulas.
There is no unbounded search
over candidate equivalents.

An empty before or after
list contributes true;
an empty current list
does the same. An empty
outer list yields false.
These are legitimate Boolean
constants in both fragments.
If a source formula simplifies
to a constant, the declared
parameter \(x\) can still
be used for its one-variable
normal form.

At a first point the before
components receive their
correct empty-past values
through the future theorem
at the last point of the
opposite order. At a last
point the ordinary future
theorem gives the empty-future
values. In a singleton
both arguments apply.
No step replaces an absent
strict witness by vacuity
of its open-gap universal.
This proves all endpoint
cases of the theorem.

For the worked identity,
reason on an arbitrary
linear order, without
the discrete hypothesis.
The left side supplies
\(y>x\) and \(z<y\)
with \(q(z)\). Classify
this same \(z\) relative
to the original point \(x\).
If \(z<x\), then \(Pq\)
holds at \(x\), and \(y\)
supplies \(F\top\).
If \(z=x\), the current
atom \(q\) and the
same \(F\top\) hold.
If \(x<z<y\), then
\(z\) witnesses
\(F(q\land F\top)\),
using \(y\) as its
later point. These
are all possible
order positions.

Conversely, a past or
current \(q\)-point
combined with any
\(y>x\) is strictly
before \(y\), so
it witnesses \(Pq\)
at that future point.
The last disjunct
supplies \(x<z<y\)
with \(q(z)\), and
again \(Pq\) holds
at \(y\). This
proves both directions
of the displayed
separated identity.
On \(0<1\) with
\(q\) only at 1,
the future \(q\)-point
has no point after
it, and neither a
past nor a current
\(q\)-point exists
at zero. Both
disjuncts are false,
as is \(F(Pq)\).
At a last point
both future existential
requirements fail.
\end{proof}
\end{annexitem}

\begin{annexitem}{annex:p9.gpss.guarded.separation.dual}{thm:p09-gpss-guarded-separation-dual}{The actual dual limit and the guarded past interface}
\begin{proof}
The explicit truth table for
\(K^-B\) in the opposite
model quantifies earlier
points there, which are
later points in the
original model. Its
witness interval becomes
\(t<u<r\). Applying the
operand reversal identity
therefore gives exactly
\(K^+_{\mathrm v}(B^{\mathrm d})\).
The converse limit calculation
is the same with the
orders exchanged. Together
with the already proved
Boolean and since/until
clauses, structural induction
establishes the extended
duality identity.

The displayed until
expression for \(K^+_{\mathrm v}\)
is obtained by negating
the existence of a
strict future neighborhood
containing no \(B\)-point.
It is true when there
is no future cut and
false when there is an
immediate successor.
These are exactly the
duals of the stipulated
past-limit endpoint values.
No nonemptiness conjunct
may be inserted silently.

A formula in \(\mathcal F\)
has only Boolean and
atomic structure outside
its until nodes, and
each past-limit occurrence
has an until ancestor.
Reversal exchanges these
with since nodes and
right-limit occurrences.
It therefore maps that
grammar into \(\mathcal P\).
Applying reversal twice
shows that it maps
onto \(\mathcal P\)
as well.

The guarded future
theorem's semantic
lemma gives tail
locality for \(\mathcal F\)
on every linear order.
By the extended
duality, a formula
of \(\mathcal P\)
has its value in
the reversed model
on the reversed
tail at its current
point. This tail
is precisely the
original initial
segment at and
before that point.
Thus every such
formula is past,
including its
allowed current
atoms.

For a before component
\(\beta\) over a
complete model,
reverse it to the
after component
\(\beta^{\mathrm d}\).
Apply the completed
guarded future
translation in
the complete
opposite model,
then reverse
the resulting
temporal formula.
Exactly the
four-step
satisfaction chain
in the preceding
proof gives
\[
             \mathcal M,t\models\beta
                 \quad\Longleftrightarrow\quad
             \mathcal M,t\models T_P^{\mathrm{com}}(\beta).
\]
Its output is
in \(\mathcal P\)
by the grammar
calculation above.
This proves the
typed past
interface at
every point.

Finally the two
limit symbols
are abbreviations
with their exact
since and until
expansions in
the body. Expand
them recursively
from their finite
syntax trees,
using fresh
first-order names
when displaying
their truth tables.
This yields an
actual \(L(S,U)\)
formula with
the same semantics.
The expansion
does not provide
a pure-operator
formula: the
guarded future
fragment can
retain since
inside its
future limit
test, and the
guarded past
fragment can
retain until
inside its
dual test.
\end{proof}
\end{annexitem}

\begin{annexitem}{annex:p9.gpss.complete.order.separation}{thm:p09-gpss-complete-order-separation}{The exact complete-order replacement of the two translators}
\begin{proof}
The temporal-to-first-order
translation and the
one-variable special
normal form hold on
all linear orders.
Thus their same
before/current/after
clause list is
available in every
complete model.
For an after
component, use
the actual guarded
future translator
\(T_F^{\mathrm{com}}\).
For a before
component, use
its fully proved
dual interface
\(T_P^{\mathrm{com}}\).
Completeness is
preserved by
reversal, so
both hypotheses
are satisfied.

Replacing the
conjuncts by
these equivalents
and retaining
their finite
Boolean combinations
proves the
output equivalent
to the original
formula in both
directions.
Its past factors
belong to
\(\mathcal P\),
its future factors
to \(\mathcal F\),
and current
Boolean literals
to both.
The previous
duality lemma
establishes their
semantic past
and future
scope on all
linear orders.
The complete-order
hypothesis belongs
to the expressive
translations, not
to this locality
calculation.

All three kinds
of empty factor
and the empty
outer disjunction
have the same
explicit Boolean
values as in
the discrete
proof. Endpoint
cases are
included by
the complete
guarded theorem
and its dual.
The two translators
and the reversal
are effective,
so the whole
algorithm is
effective.

On a discrete
complete order
one can also
reduce this
guarded output
to the two
pure fragments.
In a future
factor, traverse
Boolean and
until nodes
and replace
an encountered
entire \(K^-\)
subformula by
false. It is
evaluated after
the current
point and
has an
immediate
predecessor.
In a past
factor, traverse
Boolean and
since nodes
and replace
an encountered
entire \(K^+_{\mathrm v}\)
subformula by
false. Its
evaluation point
is before the
current one
and has an
immediate
successor.
In each case
the empty
neighbor gap
falsifies the
limit operation
independently
of its operand,
which need
not be
traversed.
The strict
guard prevents
an erroneous
replacement
of a vacuous
limit at an
extreme point.
Thus the
complete-order
comparison
is consistent
with the
discrete
pure-fragment
theorem without
claiming that
the limits
can be
removed on
all complete
orders.
\end{proof}
\end{annexitem}

\begin{annexitem}{annex:p9.gpss.pure.separation.obstruction}{thm:p09-gpss-pure-separation-obstruction}{The mixed limit formula and the two indistinguishable pure directions}
\begin{proof}
Recall the complete
models \(C_k\):
their domain is
\(\omega^k\times[0,1)\)
in lexicographic
order, and \(P\)
holds exactly at
\((\gamma,0)\).
All other proposition
letters are fixed
false. The least
point is \((0,0)\).
The existing real
model theorem proves
that the underlying
order is isomorphic
to the nonnegative
real line.

At a first point
every strict-since
formula has no
possible earlier
witness. Its
value is false
there, regardless
of its operands.
A structural
induction on a
pure \(L(S)\)
formula therefore
reduces its value
there to the
current atomic
valuation and
Boolean operations,
replacing encountered
since subformulas
by false.
The current valuation
is the same
in \(C_1,C_2\),
so all pure-since
formulas agree.

All pure-until
formulas agree
there as well.
The already proved
four-truth-set table
is independent
of \(k\) and
is Boolean closed;
its induction
applies to
arbitrary modal
nesting, with
the marked set
interpreting \(P\).
Any additional
letters are false
in both models,
so they do not
change that
induction.
Consequently every
finite Boolean
combination of
pure-since and
pure-until formulas
agrees at the
two least points.

Now expand the
mixed formula.
Its outer until
is strict
eventuality, so
it asks for
a future marked
point \(t\)
at which
\((\neg P)\mathbin{\mathsf S_s}\top\)
is false.
That since formula
is true exactly
when some earlier
cut has no
marked point
strictly between
it and \(t\).
Its negation
therefore says
that marked
points occur
arbitrarily close
below \(t\).
This is the
full semantics
of the displayed
since expansion;
no new modality
is being assumed.

In \(C_1\),
every marked
point strictly
after \((0,0)\)
has index a
positive integer
\(n\). Choose
the earlier
point \((n-1,1/2)\).
The open interval
from it to
\((n,0)\)
is unmarked.
Thus the since
formula is true,
its negation
is false, and
no future
marked point
can witness
\(\Xi\).
The initial
marked point,
where a past
universal could
be vacuous,
is excluded
by the outer
strict future
witness.

In \(C_2\),
take \(t=(\omega,0)\).
For any earlier
point \((n,r)\)
with finite
index \(n\),
the marked
point \((n+1,0)\)
lies strictly
between it
and \(t\).
These are
all points
below \(t\)
in \(C_2\).
Thus every
earlier cut
has a marked
point after
it and before
\(t\).
The since
formula is
false there,
its negation
is true, and
this future
marked point
witnesses \(\Xi\).
The two values
contradict every
proposed pure
Boolean separation
over all
complete orders.

For the stronger
full-real-line
claim, identify
each \(C_k\)
with \([0,+\infty)\)
by the proved
order isomorphism,
mapping its
least point
to zero.
Extend the
interpretation
to \(\mathbb R\)
by taking
all predicates
false at
negative points.
The initial
segments at
and before
zero are
identical in
the two models:
\(P\) is
false below
zero and true
at zero.
Pure-since
locality therefore
preserves their
agreement there.
Pure-until
locality preserves
the agreement
already proved
on the two
nonnegative
tails.

The values
of \(\Xi\)
are unchanged.
In \(C_1\)
the same
unmarked interval
before every
future marked
point still
exists. In
\(C_2\),
the original
cuts below
\((\omega,0)\)
still have
their intermediate
marked points;
any new
negative cut
has the
marked point
zero strictly
after it and
before that
witness.
Thus the
mixed formula
still separates
the two real
models, while
every pure
Boolean combination
agrees at
zero. This
proves the
claimed real
counterinstance
and the
failure of
the unqualified
pure-operator
extension.
\end{proof}
\end{annexitem}
